% Diagrama del sistema de recuperación — versión TikZ (nativa de LaTeX).
% Requiere: tikz con las librerías positioning, fit, backgrounds y shapes.geometric.
%
% NOTA: con babel en español los caracteres < y > quedan activos (shorthands) y
% rompen el parseo de las flechas -> de TikZ. Por eso el grupo siguiente los
% desactiva mientras se dibuja el diagrama; el efecto no sale del grupo.
%
% NOTA: los nodos NO usan «text width»; el ancho lo fijan su contenido y un
% «minimum width» en los cilindros. Con «text width» las palabras largas
% (Reordenamiento, Recuperación) desbordaban su caja y LaTeX avisaba de
% \overfull \hbox aunque la figura cupiera en la página. Las que no caben se
% parten a mano con «-» (Reordena-/miento, Recupe-/ración).
%
% NOTA: la fusión de rangos y el reordenamiento neuronal son ALTERNATIVAS sobre
% la salida de la recuperación, no etapas consecutivas: ambas cuelgan de la
% misma bifurcación y cada una entra a la evaluación por su propia flecha (una
% por el norte y otra por el este, para que no se empalmen).
%
% NOTA: la recuperación depende de los embeddings de documentos y de consultas
% (ambas cajas apuntan a ella) y el caché de rankings cuelga de ella con una
% doble flecha (<->): escribe ahí los rankings y los relee si ya existen; no
% es una etapa del flujo principal.
\begingroup
\ifdefined\shorthandoff\shorthandoff{<>}\fi
\begin{tikzpicture}[
  font=\scriptsize,
  x=1mm, y=1mm,
  box/.style   = {rectangle, rounded corners=2pt, draw, align=center,
                  inner sep=4pt, minimum height=9mm},
  sub/.style   = {rectangle, rounded corners=1.5pt, draw, fill=white,
                  align=center, inner sep=2pt,
                  minimum height=7mm, minimum width=24mm},
  grupo/.style = {rectangle, rounded corners=3pt, draw, fill=black!4,
                  inner sep=2mm},
  store/.style = {cylinder, shape border rotate=90, draw, align=center,
                  inner sep=3pt, minimum width=19mm,
                  minimum height=9mm, shape aspect=0.25},
  src/.style   = {rectangle, rounded corners=2pt, draw, dashed, align=center,
                  inner sep=4pt, minimum height=9mm},
  gt/.style    = {store, dashed},
  etapa/.style = {font=\tiny, text=black!70},
  every edge/.style = {draw, ->, rounded corners=2mm}
]

  %% ---------------- Columna 1: fuentes de datos ----------------
  \node[src]   (eswiki)   at (0,0)   {eswiki};
  \node[src]   (messirve) at (0,-26) {MessIRve};
  \node[etapa] at (0,17)             {Datos};

  %% ---------------- Columna 2: almacenes ----------------
  \node[store] (docs) at (25,0)   {Documentos};
  \node[store] (qrys) at (25,-26) {Consultas};

  %% ---------------- Columna 3: representación e indexación ----------------
  \node[sub] (demb) at (56,4.5)   {Vectores densos};
  \node[sub] (dinv) at (56,-4.5)  {Índice invertido};
  \node[sub] (qemb) at (56,-21.5) {Vectores densos};
  \node[sub] (qtok) at (56,-30.5) {Vectores dispersos};

  \begin{scope}[on background layer]
    \node[grupo, fit=(demb)(dinv),
          label={[etapa]above:Índices de documentos}] (dgrp) {};
    \node[grupo, fit=(qemb)(qtok),
          label={[etapa]below:Consultas codificadas}] (qgrp) {};
  \end{scope}
  \node[etapa] at (56,17) {Representación e indexación};

  %% ---------------- Columna 4: recuperación ----------------
  \node[store] (cache) at (90,0)   {Caché de\\rankings};
  \node[box]   (retr)  at (90,-26) {Recuperación\\[1pt]{\tiny densa / dispersa}};
  \node[etapa] at (90,17)          {Recuperación};

  %% ---------------- Columna 5: alternativas ----------------
  \node[box]   (fus) at (116,0)   {Fusión\\de rangos};
  \node[box]   (rer) at (116,-26) {Reordena-\\miento\\neuronal};
  \node[etapa] at (116,17)        {Alternativas};

  %% ---------------- Verdad de referencia y evaluación ----------------
  \node[gt]  (verdad) at (25,-48)  {Verdad de\\referencia};
  \node[box] (eval)   at (116,-48) {Evaluación};

  %% ---------------- Aristas: flujo principal ----------------
  \draw[->] (eswiki) -- (docs);
  \draw[->] (messirve) -- (qrys);
  \draw[->] (messirve.south) |- (verdad.west);
  \draw[->] (docs) -- (dgrp);
  \draw[->] (qrys) -- (qgrp);
  \draw[->] (dgrp.east) -- ++(4,0) -- ++(0,-13) -| ([xshift=-4mm]retr.north);
  \draw[->] (qgrp.east) -- (retr.west);
  \draw[<->] (cache.south) -- (retr.north);

  %% ---------------- Aristas: bifurcación ----------------
  \draw (retr.east) -- ++(3,0) coordinate (bif);
  \draw[->] (bif) |- (fus.west);
  \draw[->] (bif) -- (rer.west);

  %% ---------------- Aristas: evaluación ----------------
  \draw[->] (rer.south) -- (eval.north);
  \draw[->] (fus.east) -- ++(9,0) |- (eval.east);
  \draw[->, dashed] (verdad.east) -- (eval.west);

\end{tikzpicture}
\endgroup